\documentclass[12pt,a4paper]{article}

\usepackage[T1]{fontenc}
\usepackage[utf8]{inputenc}
\usepackage[polish]{babel}
\usepackage{lmodern}
\usepackage{microtype}
\usepackage{enumitem}
\usepackage{hyperref}

\title{Propozycja tematu pracy magisterskiej}
\author{}
\date{}

\begin{document}

\maketitle

\section*{Temat}

\textbf{Wybrane metody estymacji krzywej uczenia}

\section*{Literatura}

\begin{enumerate}
    \item Grigg, O., Farewell, V. T., Spiegelhalter, D. J., et al. (2003).
    The Use of Risk-Adjusted CUSUM and RSPRT Charts for Monitoring in Medical Contexts.
    \textit{Statistical Methods in Medical Research}, 12(2), 147--170.
    \href{https://doi.org/10.1177/096228020301200205}{doi:10.1177/096228020301200205}.

    \item Barnard, G. A. (1959).
    Control Charts and Stochastic Processes.
    \textit{Journal of the Royal Statistical Society. Series B (Methodological)},
    21(2), 239--271.
    \href{https://doi.org/10.1111/j.2517-6161.1959.tb00336.x}
    {doi:10.1111/j.2517-6161.1959.tb00336.x}.

    \item Basseville, M., Nikiforov, I. V. (1993).
    \textit{Detection of Abrupt Changes: Theory and Application}.
    Englewood Cliffs, NJ: Prentice-Hall.

    \item Cook, J. A., Ramsay, C. R., Fayers, P. (2004).
    Statistical Evaluation of Learning Curve Effects in Surgical Trials.
    \textit{Clinical Trials}, 1(5), 421--427.

    \item Park, K., Kwon, J., Yoo, E., Lee, S., Song, J., Pyeon, S. (2024).
    Application of CUSUM Analysis in Assessing Learning Curves in Robot-Assisted
    Sacrocolpopexy Performed by Experienced Gynecologist.
    \textit{BMC Surgery}, 24.
\end{enumerate}

\end{document}
